% ===========================================================================
\section{The Parallel-Table Trap in the IPCC Defaults}
\label{sec:factors}
% ===========================================================================

The 2006 Guidelines \citep{ipcc2006} publish stationary and mobile combustion factors not as
one table per fuel but as several parallel tables, one per source category. Volume~2
Chapter~2 gives four for stationary combustion---energy industries, manufacturing,
commercial, and residential---and Chapter~3 gives separate tables for road, rail and off-road
mobile sources. The same fuel appears in all of them.

Across those tables the \COtwo{} factor for a given fuel is \emph{identical}, because \COtwo{}
depends only on the carbon in the fuel and on complete oxidation. The \CHfour{} and \NtwoO{}
factors are not identical, and in some cases differ by nearly an order of magnitude, because
they depend on how completely and at what temperature the fuel burns.

This creates a failure mode with an unusually bad signature. A compiler who selects the wrong
table gets \COtwo{} exactly right---and \COtwo{} is typically \SIrange{95}{99}{\percent} of a
fossil combustion category, so every obvious check passes. The total looks correct. The
comparison against an external \COtwo{} dataset agrees. Only the trace gases are wrong, and
they can be wrong by a factor large enough to move the affected category by
\SI{8}{\percent}---which is smaller than most published uncertainty bands and therefore
invisible to any validation that works at the level of the total.

\textbf{How much of the total moves depends on how \COtwo{}-dominated the category is, and an
earlier edition of this paper overstated it.} The railway and mixed off-road cases swap
\NtwoO{} between \SI{28.6}{} and \SI{3.9}{\kilo\gram\per\tera\joule} and move their own
category by \SI{8.0}{\percent}. The power and residential cases swap \CHfour{} on gas streams
whose totals are almost entirely \COtwo{}, and move the sector by \SI{0.08}{} and
\SI{0.12}{\percent} respectively---the \CHfour{} row moves several-fold, the total barely at
all. This paper previously reported all four as \SIrange{8}{9}{\percent} sector errors. The
magnitudes are now computed from the engine's activity and the two published factors by
\texttt{tests/test\_factor\_table\_trap.py}, so the claim cannot be restated without the
arithmetic agreeing.

The finding is unharmed by the correction, and is arguably sharper for it. What makes these
errors dangerous is not how much of a total they move but that they move NONE of the
\COtwo{}: every external comparator this inventory has is \COtwo{}-dominated, so all four are
invisible to all of them regardless of size.

We encountered this trap four times in ten sectors. Table~\ref{tab:factortrap} gives all
four instances with the quantified consequence of the plausible error. Each is pinned by a
test that fails if the wrong table's value is ever substituted, and each sector's module
carries the warning in its own documentation with a pointer to its mirror image.

\begin{table*}[htbp]
\centering
\small
\begin{tabular}{@{}p{0.15\linewidth}p{0.13\linewidth}p{0.13\linewidth}p{0.22\linewidth}p{0.27\linewidth}@{}}
\toprule
\textbf{Sector} & \textbf{Correct table} & \textbf{Plausible error} & \textbf{Factor
difference} & \textbf{Consequence of the error} \\
\midrule
Electricity generation (1A1a) &
Vol~2 Ch~2 Table~2.2, Energy Industries &
Table~2.5, Residential &
Natural gas \CHfour{} \SI{1}{\kilo\gram\per\tera\joule} vs \num{5} &
\COtwo{} unchanged; \CHfour{} inflated fivefold. A utility boiler burns far more completely
than a domestic appliance \\[3pt]

Residential fuel (1A4b) &
Vol~2 Ch~2 Table~2.5, Residential &
Table~2.2, Energy Industries &
Natural gas \CHfour{} \SI{5}{\kilo\gram\per\tera\joule} vs \num{1}; kerosene \num{10} vs
\num{3} &
The mirror image: \CHfour{} understated by \SI{80}{\percent} on a quarter of household
energy input \\[3pt]

Railways (1A3c) &
Vol~2 Ch~3 Table~3.4.1, Railways &
Table~3.2.2, Road diesel &
\NtwoO{} \SI{28.6}{\kilo\gram\per\tera\joule} vs \num{3.9}---\num{7.3} times higher &
Railway \NtwoO{} understated by \SI{86}{\percent} and the sector total by about
\SI{9}{\percent}. Large slow-speed locomotive engines run lean at high load, the
\NtwoO{}-forming regime \\[3pt]

Off-road machinery (1A4c) &
Vol~2 Ch~3 Table~3.3.1, Off-road &
Table~3.2.2, Road diesel &
\NtwoO{} \SI{28.6}{\kilo\gram\per\tera\joule} vs \num{3.9} &
The category would report about \SI{8}{\percent} less. Identical to the rail value because
both derive from EEA~2005 and neither engine class generally carries a catalyst \\
\bottomrule
\end{tabular}
\caption{Four instances of the parallel-table trap in this inventory. In every case the
plausible wrong table leaves \COtwo{} exactly correct---so the sector's dominant gas, its
external benchmark comparison and its mass balance all still agree---while the trace gases
move. How far the total moves depends on how \COtwo{}-dominated it is: the railway and
off-road swaps move their own \emph{category} by \SI{8.0}{\percent}, the power and residential
swaps move the \CHfour{} row several-fold and the \emph{sector total} by \SI{0.06}{} and
\SI{0.12}{\percent}. Each magnitude is computed from the engine by a regression test rather
than quoted.}
\label{tab:factortrap}
\end{table*}

\subsection{Why this is a compiler-facing finding rather than a coding note}

Three properties make this worth reporting rather than fixing quietly.

It is \emph{systematic}: it arose in four of the six fuel-burning sectors, and in each case
the wrong table was the one a reasonable person would reach for. The power sector's fuel is
delivered to buildings; the residential sector's largest fossil fuel is natural gas, whose
most familiar factor is the utility one; rail and off-road both burn ordinary road diesel
bought from the same national supply.

It is \emph{invisible to validation}. We compare each sector against EDGAR, FAOSTAT, Climate
TRACE and, where available, Pakistan's own BUR-1. Every one of those comparators is dominated
by \COtwo{} for fossil categories. A sector using the wrong trace-gas table passes all of
them. The only thing that catches it is reading the guidelines carefully enough to notice
that there is more than one table, which is precisely the step under time pressure.

It is \emph{directional}. The errors do not average out. Three of the four instances would
understate the trace gases, because the road and energy-industry tables carry the lower
\CHfour{} and \NtwoO{} values and are the more familiar tables. An inventory community making
this error at random rates would be biased low on non-\COtwo{} combustion emissions.

Two related decisions in the same family are worth recording because they run the other way.
First, we apply \emph{no separate oxidation factor} anywhere in fossil combustion: the 2006
Guidelines assume complete oxidation and fold it into the emission factor, with Table~1.4
carrying an explicit oxidation column of \num{1.0} for every fuel, so applying the
1996-vintage \num{0.98} on top would double-discount. Second, and deliberately contrasting,
the waste sector \emph{does} carry an oxidation factor of \num{0.71}, because open burning
genuinely is incomplete. Both are pinned by tests, so that neither can be ``corrected'' into
the other.

Finally, the power sector's activity is handled on an energy basis rather than a mass basis
for a reason from the same family. IPCC's bituminous and lignite \COtwo{} factors differ by
only \SI{6.8}{\percent} per terajoule but by a factor of two per tonne of coal, because their
net calorific values differ by \SI{54}{\percent} and lignite's own published range spans
\SIrange{5.50}{21.6}{\tera\joule\per\giga\gram}. Working in energy units makes that
uncertainty disappear rather than requiring an assumption about which coal Pakistan burns.
